% Created 2023-02-13 lun 17:29
% Intended LaTeX compiler: pdflatex
\documentclass[11pt]{article}
\usepackage[utf8]{inputenc}
\usepackage[T1]{fontenc}
\usepackage{graphicx}
\usepackage{longtable}
\usepackage{wrapfig}
\usepackage{rotating}
\usepackage[normalem]{ulem}
\usepackage{amsmath}
\usepackage{amssymb}
\usepackage{capt-of}
\usepackage{hyperref}
\usepackage{minted}

\usepackage{parskip}
\usepackage[utf8]{inputenc} %% For unicode chars
\usepackage{placeins}
\usepackage[margin=2.50cm]{geometry}
\usepackage[T1]{fontenc}
\usepackage{mathpazo}
\linespread{1.05}
\usepackage[scaled]{helvet}
\usepackage{courier}
\hypersetup{colorlinks=true,linkcolor=blue}
\RequirePackage{fancyvrb}
\DefineVerbatimEnvironment{verbatim}{Verbatim}{fontsize=\small,formatcom = {\color[rgb]{0.5,0,0}}}
\usepackage{mdframed}
\BeforeBeginEnvironment{minted}{\begin{mdframed}}
\AfterEndEnvironment{minted}{\end{mdframed}}
\date{}
\title{Redes Neuronales: Backpropagation}
\hypersetup{
 pdfauthor={},
 pdftitle={Redes Neuronales: Backpropagation},
 pdfkeywords={},
 pdfsubject={},
 pdfcreator={Emacs 27.1 (Org mode 9.6)}, 
 pdflang={Spanish}}
\begin{document}

\maketitle
\setlength\parindent{10pt}


\section*{Estudio de la implementación de Backpropagation}
\label{sec:org31f8a86}

Comprender el algoritmo: Es importante comprender el funcionamiento
básico del algoritmo de Backpropagation antes de compararlo con el
código. Esto incluye cómo se realiza la propagación hacia atrás y cómo
se ajustan los pesos en la red neuronal.

Estudia la implementación en Python del algoritmo de Backprogagation y
analiza los siguientes puntos:


\subsection*{Tareas}
\label{sec:org45244c6}
\begin{itemize}
\item \textbf{Correspondencia con el código}. \textbf{Identifica} los pasos del
algoritmo en el código y asegúrate de que \textbf{el código refleje
correctamente el algoritmo} y que se haya implementado de manera
adecuada.

\item \textbf{Eficiencia y escalabilidad}: La eficiencia y escalabilidad del
código son factores importantes a considerar en la
comparativa. \textbf{Evalúa el tiempo de ejecución y la capacidad de la
implementación para manejar grandes conjuntos de datos}.

\item \textbf{Límites y desventajas}: Evalúa los límites y desventajas tanto
del algoritmo como del código. Esto incluye factores como la
sensibilidad a la tasa de aprendizaje, la necesidad de una gran
cantidad de datos para el entrenamiento y la posibilidad de
quedarse atrapado en mínimos locales.

\item \textbf{Mejoras y optimizaciones}: Mira si puedes hacer cualquier mejora u
optimización en el código, como el uso de técnicas de
regularización..
\end{itemize}
\end{document}